\chapter{Paradigma Algorítmico e Modelagem}
\label{cap:paradigma}

% TODO(P2): Apresentar os algoritmos completos em pseudocódigo utilizando o pacote algorithm2e
% e os diagramas conceituais em TikZ obrigatórios conforme o enunciado:
%
% 1. Sequência de Fibonacci:
%    - Algoritmo Fib-Naive (recursivo exponencial).
%    - Algoritmo Memoized-Fib (com tabela de cache).
%    - Algoritmo Bottom-Up-Fib (iterativo com vetor).
%
% 2. Corte de Hastes (Cormen 15.1):
%    - Algoritmo Cut-Rod (ingênuo).
%    - Algoritmo Memoized-Cut-Rod e auxiliar Memoized-Cut-Rod-Aux.
%    - Algoritmo Bottom-Up-Cut-Rod.
%    - Algoritmo Extended-Bottom-Up-Cut-Rod (retornando r e s).
%    - Algoritmo Print-Cut-Rod-Solution (reconstrução iterativa).
%
% 3. Multiplicação em Cadeia de Matrizes (Cormen 15.2):
%    - Algoritmo Matrix-Chain-Order (preenchimento por diagonais com laço triplo).
%    - Algoritmo Print-Optimal-Parens (reconstrução recursiva da parentização).
%
% 4. Subsequência Comum Máxima - LCS (Cormen 15.4):
%    - Algoritmo LCS-Length (tabelas c e b).
%    - Algoritmo Print-LCS (traceback da subsequência ótima).
%
% 5. Diagramas Obrigatórios em TikZ (em relatorio/tikz/):
%    - Diagrama (a): Árvore de recursão do Cut-Rod para n = 4 com subproblemas repetidos coloridos.
%    - Diagrama (b): Matriz de programação dinâmica do LCS preenchida com setas direcionais ('diag', 'up', 'left').
%    - Diagrama (c): Grafo direcionado acíclico (DAG) de dependências dos subproblemas.
%    - Diagrama (d): Tabelas m e s completas para o exemplo de 6 matrizes do Cormen (dims = [30, 35, 15, 5, 10, 20, 25]).
%    - Diagrama (e - opcional): Estado da pilha de execução e cache durante a memoização.

% [Texto, pseudocódigos e diagramas TikZ a cargo de P2 - ver docs/tarefas/P2_relatorio_secoes_3_4.md]
